\documentclass[11pt]{article}

% ── Packages ────────────────────────────────────────────────
\usepackage[margin=1in]{geometry}
\usepackage{amsmath, amssymb, amsthm}
\usepackage{booktabs}
\usepackage{graphicx}
\usepackage{hyperref}
\usepackage{natbib}
\usepackage{setspace}
\usepackage{microtype}
\usepackage{enumitem}
\usepackage{caption}
\usepackage{subcaption}
\usepackage{xcolor}
\usepackage{float}
% code packages
\usepackage{listings}
\usepackage{inconsolata}
\lstset{basicstyle=\ttfamily\footnotesize,breaklines=true}
\usepackage{csvsimple}
\usepackage{booktabs}

% ── Hyperref setup ──────────────────────────────────────────
\hypersetup{
    colorlinks = true,
    linkcolor  = black,
    citecolor  = black,
    urlcolor   = blue
}

% ── Theorem environments ────────────────────────────────────
\newtheorem{theorem}{Theorem}
\newtheorem{proposition}{Proposition}
\newtheorem{lemma}{Lemma}
\newtheorem{corollary}{Corollary}
\theoremstyle{definition}
\newtheorem{definition}{Definition}
\newtheorem{assumption}{Assumption}
\theoremstyle{remark}
\newtheorem{remark}{Remark}

% ── Spacing ─────────────────────────────────────────────────
\onehalfspacing

% ── Title block ─────────────────────────────────────────────
\title{Accounting for Love of Variety in Consumer Demand Estimation}
\author{Tate Mason}
\date{\today}

% ============================================================
\begin{document}
\maketitle
\thispagestyle{empty}

%============================================================%
\section{Introduction}
%============================================================%

In the baseline logit model, a consumer is perceived as having a fixed preference for each product, and the utility from consuming a product is independent of past consumption. Further, it is assumed that a new good is always a utility boosting event, by the error term's property of being identically and independently distributed. However, in many markets, consumers may have a love of variety, or a desire to consume different products over time. This can be due to factors such as boredom, satiation, or a preference for novelty. In such cases, the utility from consuming a product may depend on past consumption, and the introduction of a new good may not always lead to an increase in utility.

% contribute dynamics and supply side effects to LOV literature

I find that...

The rest of the paper is organized as follows: section two discusses the related
literature and section three describes the data used and some descriptive
statistics. Section four showcases the model and environment and section five
presents results. Section six goes into counterfactual analysis, while section 7
concludes.

%============================================================%
\section{Related Literature}
%============================================================%

    The variety-seeking and state dependence literature has been primarily focused on understanding the dynamics of consumer preferences and demand over time. From \cite{lattin_model_1987}, who developed an empirical structural model of state dependence in consumer choice, to \citet{erdem_decision-making_1996}, who iterated the model to allow for consumer learning. \cite{chintagunta_variety_1999} is the most closely related paper, as he develops a model of variety-seeking behavior which is both tractible and flexible, and uses Nielsen data to estimate the model. However, this paper only looks at demand, leaving me a niche to explore the interplay between variety-seeking and strategic advertising.

    % Dube (2004), McAllister (1984), Train (2010), Chintgunta, Dube, Hao (2006)


%============================================================%
\section{Data}
%============================================================%

    \subsection{Data}

    My data comes from Nielsen's HomeScan and MarketScan data. These two main
    datasets give me access to household purchase decisions as well as product
    characteristics. I am looking at the years between 2014-2016 in markets
    Atlanta, Chicago, San Diego, and Denver. Households are single-person to
    avoid any contamination of shopping for multiple people like children or
    spouses, allowing me a cleaner sample of shopping trips. Stores have to be valid Nielsen stores, meaning they appear in the reatil code data. This is necessary to observe prices the products were bought at. 

    The product space I am looking at is yogurt. Yogurt is good for this sort of
    analysis due to its relative lack of storability as well as its variation in
    attributes like flavor, nutrient contents, and types like Greek or Skyr. I
    have limited yogurt sizes to two kinds: cups and tubs. Cups of yogurt are
    between five and seven ounces while tubs are approximately 32 ounces. Cups
    are single serving products and come in more varied flavors. Tubs, on the
    other hand, tend to be limited in flavors but hold five to seven servings of
    yogurt. These two types provide interesting anaylsis of what type of person
    buys which size.

    \ref{tab:stat_tab} shows that there are n households, with a in Atlanta, b
    in Chicago, c in Denver, and d in San Diego. The mean income is \$17,000 and
    the mean age is n years old. The sample is mostly white at n\%. Households
    go on an average of n shopping trips, indexed by a unique trip code. Per
    trip, they buy an average of n units of yogurt. If we break out by type, n
    cups and/or n tubs. 

    I define the outside option to be anytime yogurt isn't purchased during a trip. This occurred
    x\%, on average.

    \ref{fig:flav_heatmap} shows that people who enjoy plain yogurt stay with their flavor longer, staying with plain for about 6 periods as opposed to about 2 for flavored people. This data fact implies that people do differ in their preferences for variety along the flavor dimension. Further along these lines, we see that \ref{fig:var_timeseries} shows ... . This is indicative of ... .
%============================================================%
\section{Model}
%============================================================%

    \subsection{Environment}
    % customers, trips, occasions, menu in store

    There are a continuum of $i \in \mathcal{I}$ consumers who exist for
    $\mathcal{T}$ periods. Within a given period, a household can go for
    a shopping trip $t\in\mathcal{T}$ which is comprised of choice occasions
    $o\in\mathcal{O}_t$. Each choice occasion allows for the purchase of a
    singular good. 


    \subsection{Consumer's Problem}
    The consumer's problem is the following. In each period $t$, the consumer $i$ chooses a product $j$ from the menu of available products $J_t$. Each product $j\in J_t$ is a bundle of $k\in K$ characteristics such that $X_{jt}^k \in \mathcal{X}_{jt}$. The consumer's utility from choosing product $j$ at time $t$ is given by:


    \begin{alignat*}{1}
      U_{ijt} &= \delta D_{ht} + \sum_1^k\beta^k_i X^k_{jt} + \gamma_i \log(1 + \Xi_{ijt}) 
                  + \alpha_i p_{jt} + \sigma\hat{\eta} + \xi_j + \varepsilon_{ijt} \\
    \text{subject to} \\
      \Xi_{ijt} &= \sqrt{\sum_{k}(X^k_{jt} - \theta^k_{ijt})^2} \\
      \theta^k_{ijt} &= \lambda\frac{1}{t-2}\sum^{t-2}_{s} X^{*k}_{js} +
      (1-\lambda)X^{*k}_{jt-1} \\
      \hat{\eta} = p_jt - \bar{p} - \tau
      \gamma_i &\sim \mathcal{N}(0, 1)\\ \varepsilon_{ijt} &\sim \text{T1EV}(0,1)

    \end{alignat*}

    $\beta^k_i$ is a consumer's utility from characteristic $X^k_{jt}$, $\gamma_i$ is their preference for variety, $\Xi_{ijt}$ is the difference between the current characteristic and the previous choices, $\theta^k_{it}$. $\theta^k_{it}$ is a function of the consumers' past purchases, where $\lambda$ weights chosen products in the past, $X^{*k}_{js}$ and the most recently chosen product, $X_{jt-1}^{*k}$, allowing for more recent purchases to play more of a role in the current period's decision. This could take many forms, but in this paper will be the mean characteristic of past purchases. $\alpha_i$ is a hetergeneous price sensitivity to $p_{jt}$. $\xi_j$ is unobserved quality in product $j$, and $\varepsilon_{ijt}$ is the Gumbel shock which allows for a closed form logit.

    % Control function to correct for price endogeneity, Dube arrival rate.


    The choice probability of consumer $i$ choosing product $j$ in time $t$ is as follows:
    \begin{alignat*}{1}
      P_{ijt} = \frac{\exp{U_{ijt}}}{\sum_k\exp{U_{ikt}}}
    \end{alignat*}

%============================================================%
\section{Results}
%============================================================%

% base tables
% marginal effects
% dollarized

%============================================================%
\section{Counterfactuals}
%============================================================%


    % Supply side aspects

    % Counterfactual sales on products
    
    % $.20 sale leads to x% change in switching likelihood

    % product intro (new flavor category ?)


%============================================================%
\section{Conclusion}
%============================================================%

In closing, this paper develops a dynamic discrete choice model of consumer demand that incorporates a love of variety and strategic advertising by firms. The model is estimated using scanner data, and counterfactual analyses are conducted to evaluate the effects of different market structures on consumer welfare and market outcomes. The paper contributes to the literature on  advertising and consumer behavior, by providing empirical evidence on the effects of strategic advertising in a dynamic demand setting. The findings of this paper have important implications for policymakers and firms, as they highlight the potential benefits and drawbacks of targeted advertising in markets when accounting for consumers' preference for variety.

%============================================================%
\newpage

\bibliography{lov_lit.bib}

\end{document}

